% theme: dot_product_angles_and_perpendicularity
\MajorQuestionLesson{内積となす角・垂直}
\SetRowSpace{5mm}

\TypeQuestionSingle{次の2つのベクトルのなす角 $\theta$ を求めなさい。}{angle_between_vectors}
\QQRow{$(3,4),\ (4,-3)$}{$90^\circ$}{$(1,2),\ (3,1)$}{$45^\circ$}
\QQRow{$(1,0),\ (-1,\sqrt3)$}{$120^\circ$}{$(\sqrt3,1),\ (3,0)$}{$30^\circ$}
\QQRow{$(2,-2),\ (-1,0)$}{$135^\circ$}{$(-2,2\sqrt3),\ (1,\sqrt3)$}{$60^\circ$}

\TypeQuestionDouble{次の2つのベクトルが垂直になるように、実数 $k$ の値を求めなさい。}{perpendicular_parameter}
\QQRow{$(2,3),\ (k,-2)$}{$k=3$}{$(k,4),\ (2,3)$}{$k=-6$}
\QQRow{$(k-1,2),\ (3,k)$}{$k=\dfrac35$}{$(1,k+2),\ (2k,1)$}{$k=-\dfrac23$}

\TypeQuestionSingle{次のベクトルに垂直な単位ベクトルをすべて求めなさい。}{perpendicular_unit_vector}
\QQRowThree{$(3,4)$}{$\pm\left(-\dfrac45,\dfrac35\right)$}{$(-5,12)$}{$\pm\left(-\dfrac{12}{13},-\dfrac5{13}\right)$}{$(1,-1)$}{$\pm\left(\dfrac{\sqrt2}{2},\dfrac{\sqrt2}{2}\right)$}
\QQRowThree{$(0,7)$}{$\pm(1,0)$}{$(2,\sqrt5)$}{$\pm\left(-\dfrac{\sqrt5}{3},\dfrac23\right)$}{$(-8,-6)$}{$\pm\left(\dfrac35,-\dfrac45\right)$}

\TypeQuestionDouble{$|\vec{a}|,|\vec{b}|$ と $\vec{a}\cdot\vec{b}$ が次のように与えられている。2つのベクトルのなす角を求めなさい。}{angle_from_dot_product}
\QQRow{$|\vec{a}|=2,\ |\vec{b}|=3,\ \vec{a}\cdot\vec{b}=3$}{$60^\circ$}{$|\vec{a}|=4,\ |\vec{b}|=5,\ \vec{a}\cdot\vec{b}=-10$}{$120^\circ$}
\QQRow{$|\vec{a}|=3,\ |\vec{b}|=6,\ \vec{a}\cdot\vec{b}=0$}{$90^\circ$}{$|\vec{a}|=\sqrt2,\ |\vec{b}|=\sqrt6,\ \vec{a}\cdot\vec{b}=\sqrt3$}{$60^\circ$}

\TypeQuestionSingle{次の条件を満たす実数 $k$ を求めなさい。}{specified_angle_parameter}
\QQRow{$(1,0)$ と $(k,1)$ のなす角が $45^\circ$}{$k=1$}{$(1,\sqrt3)$ と $(k,1)$ のなす角が $60^\circ$}{$k=-\dfrac{\sqrt3}{3}$}
\QQRow{$(2,1)$ と $(1,k)$ のなす角が $90^\circ$}{$k=-2$}{$(1,-1)$ と $(k,2)$ のなす角が $135^\circ$}{$k=0$}

\LessonPageBreak
\SetRowSpace{5mm}

\TypeQuestionDouble{次の2つのベクトルのなす角が鋭角、直角、鈍角のいずれであるかを答えなさい。}{classify_angle_by_dot_product}
\QQRow{$(2,1),\ (-1,3)$}{鋭角}{$(3,-2),\ (2,3)$}{直角}
\QQRow{$(-4,1),\ (1,2)$}{鈍角}{$(-2,-3),\ (-4,1)$}{鋭角}

\TypeQuestionDouble{次の3点について、指定された角を求めなさい。}{angle_from_three_points}
\QQ{$\mathrm{A}(3,1),\ \mathrm{B}(2,1),\ \mathrm{C}(0,3)$ の $\angle\mathrm{ABC}$}{$135^\circ$}
\QQ{$\mathrm{A}(3,-2),\ \mathrm{B}(1,-2),\ \mathrm{C}(3,0)$ の $\angle\mathrm{ABC}$}{$45^\circ$}
\QQ{$\mathrm{A}(-3,-2),\ \mathrm{B}(-2,-1),\ \mathrm{C}(0,1)$ の $\angle\mathrm{ABC}$}{$180^\circ$}
\QQ{$\mathrm{A}(0,3),\ \mathrm{B}(-1,2),\ \mathrm{C}(-2,3)$ の $\angle\mathrm{ABC}$}{$90^\circ$}

\TypeQuestionSingle{次の問いに答えなさい。}{dot_product_applications}
\QQ{$\vec{a}=(2,-1)$、$\vec{b}=(3,4)$ とする。$\vec{p}=\vec{a}+t\vec{b}$ が $\vec{a}$ と垂直になるように、実数 $t$ の値を求めなさい。}{$t=-\dfrac52$}
\QQ{$|\vec{a}|=3$、$|\vec{b}|=\sqrt6$ であり、$2\vec{a}+3\vec{b}$ と $\vec{a}-\vec{b}$ が垂直である。$\vec{a}$ と $\vec{b}$ のなす角を求めなさい。}{$90^\circ$}
\QQ{$|\vec{a}|=3$、$|\vec{b}|=2$、$\vec{a}\cdot\vec{b}=-2$ とする。$|\vec{a}+t\vec{b}|$ の最小値と、そのときの実数 $t$ の値を求めなさい。}{最小値 $2\sqrt2$、$t=\dfrac12$}

\TypeQuestionDouble{次の3点を頂点とする三角形の面積を求めなさい。}{triangle_area_from_coordinates}
\QQ{$\mathrm{A}(1,-2),\ \mathrm{B}(5,1),\ \mathrm{C}(-1,4)$}{$15$}
\QQ{$\mathrm{A}(1,2),\ \mathrm{B}(7,2),\ \mathrm{C}(4,6)$}{$12$}
\QQ{$\mathrm{A}(-3,1),\ \mathrm{B}(2,-2),\ \mathrm{C}(4,5)$}{$\dfrac{41}{2}$}
\QQ{$\mathrm{A}(-2,-1),\ \mathrm{B}(-2,5),\ \mathrm{C}(4,2)$}{$18$}
